\begin{minipage}[t][][t]{0.50\linewidth}
    Représenter le vecteur champ magnétique $\vec{B}$ créé par le conducteur aux
    points $M_{1}$ et $M_{2}$.~\textbf{(0,5)}

    En déduire le sens du courant $I$.~\textbf{(0,5)}
\end{minipage}
\hfill
\begin{minipage}[t][][t]{0.48\linewidth}
    \begin{figure}[H]
        \centering
        \begin{tikzpicture}
            \def\d{5}
            \def\ar{1.6}
            % Aiguille
            \node[diamond,inner sep=0pt,minimum width=7mm,
                  minimum height=3.4mm] (aa) at (279:\ar) {};
            \path[draw,fill=white,rounded corners=0.4pt]
                (aa.north) -- (aa.south) -- (aa.east) -- cycle;
            \path[draw,fill=red,rounded corners=0.4pt]
                (aa.north) -- (aa.south) -- (aa.west) -- cycle;
            \node[below,inner sep=4pt,font=\scriptsize] at (aa.east) {S};
            \node[below,inner sep=4pt,font=\scriptsize] at (aa.west) {N};
            \draw[] (0,0) circle (\ar);
            \draw[<-,shorten <=2pt] (180:\ar) -- ++(-1,0)
                node[left,align=center,font=\scriptsize] {Ligne de\\ champ};
            % M
            \node[circle,fill,inner sep=1pt,
                label={[above]:$M_{1}$}] (M) at (aa.center) {};
            \node[circle,fill,inner sep=1pt,
                label={[above right=-1pt,inner sep=0pt]:$M_{2}$}] at (45:\ar) {};
            % fil
            \node[circle,fill=gray,inner sep=1pt,
                label={[above,inner sep=4pt,font=\scriptsize,
                align=center]:{fill\\ rectiligne}}] at (0,0) {};
        \end{tikzpicture}
    \end{figure}
\end{minipage}


	


